\subsection{Anforderungen für Web-, Cloud- und Desktop-Anwendungen}
\label{subsec:anforderungen-plattformen}

Alle Varianten sollen Frontend-Basis, Struktur, Konfigurationslogik, Tests und
Dokumentation teilen. Browserprojekte benötigen keine Desktop-Schicht und nur
bei Bedarf ein Backend. Cloudvarianten erfordern API, Deployment-Grenzen sowie
Health- und Readiness-Prüfung. Desktopvarianten verwenden dasselbe Frontend in
einer nativen Laufzeit; kombinierte Varianten ergänzen Backend und
gegebenenfalls Browserzugriff.

\begin{table}[H]
  \centering
  \scriptsize
  \renewcommand{\arraystretch}{1.08}
  \caption{Anforderungsmatrix der unterstützten Projektarten}
  \label{tab:plattform-anforderungen}
  \begin{tabularx}{\textwidth}{@{}p{0.24\textwidth}>{\centering\arraybackslash}X>{\centering\arraybackslash}X>{\centering\arraybackslash}X>{\centering\arraybackslash}X@{}}
    \toprule
    \textbf{Anforderung} & \textbf{Web} & \textbf{Cloud} & \textbf{Desktop lokal} & \textbf{Desktop + Cloud} \\
    \midrule
    Frontend & erforderlich & erforderlich & erforderlich & erforderlich \\
    Backend / API & optional & erforderlich & nicht grundsätzlich & erforderlich \\
    Desktop-Laufzeit & nein & nein & erforderlich & erforderlich \\
    Deployment-Grenze & optional & erforderlich & nein & erforderlich \\
    Persistenz & optional & optional & produktspezifisch & optional \\
    \bottomrule
  \end{tabularx}
  \smallskip
  {\footnotesize Quelle: Eigene Darstellung auf Grundlage von
  \textcite{kleiveist2026profiles}.\par}
\end{table}

PostgreSQL ist nur eine optionale Capability für backendfähige Profile, kein
eigenes Plattformprofil \parencite{kleiveist2026profiles}. Öffentliche
Client-Konfiguration und serverseitige Werte einschließlich Geheimnissen sind
zu trennen; eine bestimmte Cloud-Plattform wird nicht vorausgesetzt.
\beginnerinput{workspace/statement/200_anforderungen/230}
